\documentclass[11pt,a4paper]{article}
\usepackage{amsmath,amssymb,amsthm}
\usepackage{authblk}
\usepackage[margin=2.5cm]{geometry}
\usepackage{hyperref}
\usepackage{booktabs}
\usepackage{enumerate}
\usepackage{xcolor}

% Theorem environments
\newtheorem{theorem}{Theorem}
\newtheorem{proposition}[theorem]{Proposition}
\newtheorem{corollary}[theorem]{Corollary}
\newtheorem{conjecture}[theorem]{Conjecture}
\newtheorem{lemma}[theorem]{Lemma}
\theoremstyle{definition}
\newtheorem{definition}[theorem]{Definition}
\newtheorem{observation}[theorem]{Observation}
\newtheorem{remark}[theorem]{Remark}

% Notation
\newcommand{\ZZ}{\mathbb{Z}}
\newcommand{\QQ}{\mathbb{Q}}
\newcommand{\CC}{\mathbb{C}}
\newcommand{\HH}{\mathbb{H}}
\newcommand{\Zfive}{\ZZ/5}
\newcommand{\PSL}{\mathrm{PSL}(2,\CC)}
\newcommand{\vol}{\mathrm{vol}}
\newcommand{\Mn}{M_n}
\newcommand{\Mcusp}{m003_{\mathrm{cusp}}}
\newcommand{\MPMNS}{M_{\mathrm{PMNS}}}
\newcommand{\vp}{\varphi}

\begin{document}

\title{A Quadratic Cover Prime Formula for a Farey Tower\\
of Arithmetic Hyperbolic 3-Manifolds}

\author{Marvin L. Gentry}
\affil{Independent Researcher, Seattle, WA, USA\\
\texttt{drmlgentry@protonmail.com}\quad
ORCID: 0009-0006-4550-2663}
\date{\today}

\maketitle

\begin{abstract}
Let $\Mcusp$ denote the cusped hyperbolic 3-manifold with
SnapPy identifier \texttt{m003}, which has trace field
$\QQ(\sqrt{-3})$ and cusp shape $\tau = e^{i\pi/3}$ (the primitive
sixth root of unity $\omega$).
We study the infinite Farey tower of Dehn fillings
$\Mn = m003(-n,\,2n+1)$, $n \geq 1$.
Each $\Mn$ is a closed orientable hyperbolic 3-manifold with
$H_1(\Mn;\ZZ) = \Zfive$; the first element $M_1$ is the Meyerhoff
manifold (SnapPy \texttt{OrientableClosedCensus[1]}).

We establish three results.
First, the surgery formula
$|H_1(m003(p,q))| = 5|2p+q|$
holds for all hyperbolic Dehn fillings of $\Mcusp$,
giving $|H_1(\Mn)| = 5$ for every $n$ (since $2(-n)+(2n+1) = 1$).
Second, the cusp shape of $\Mcusp$ is exactly $\tau = \omega$
(verified to $< 10^{-15}$ residual), reflecting the equilateral
Eisenstein symmetry of the trace field $\QQ(\sqrt{-3})$.
Third, we discover and verify computationally (for $n = 1,\ldots,6$)
the following:

\medskip
\emph{The unique degree-5 algebraic covering space of $\Mn$ has
$H_1 \cong (\Zfive)^2 \oplus \ZZ/p_n$ where
\[
  p_n = 5n(n+1)+1.
\]}

The sequence begins $11, 31, 61, 101, 151, 211$ (all prime for $n \leq 7$).
The prime $p_1 = 11$ is the Lucas number $L_5$; no other $p_n$ is Lucas.
This establishes $M_1 = \MPMNS$ as the unique \emph{Lucas-prime} element
of the tower, providing a canonical arithmetic characterization of the
Meyerhoff manifold within its Farey family.

The volumes $\vol(\Mn)$ increase monotonically to $\vol(\Mcusp) = 2.02988$
with the Neumann--Zagier asymptotic $\vol(\Mcusp) - \vol(\Mn) \sim C/L_n^2$
where $L_n^2 = 5n^2+4n+1$ and $C \approx 14.16$.
\end{abstract}

\tableofcontents

%------------------------------------------------------------
\section{Introduction}
\label{sec:intro}
%------------------------------------------------------------

Compact hyperbolic 3-manifolds obtained by Dehn surgery on a
one-cusped hyperbolic 3-manifold form natural infinite families
parametrised by the Dehn surgery slope $(p,q) \in \ZZ^2$,
$\gcd(|p|,q)=1$.
By the Thurston--Jorgensen theorem~\cite{Thurston78}, the volumes
of these fillings converge to the volume of the cusped manifold,
and for all but finitely many slopes the filled manifold is hyperbolic.

In this paper we study the family of Dehn fillings of the cusped
manifold $\Mcusp$ (SnapPy identifier \texttt{m003}) along the
Farey ray of slopes $(-n,\,2n+1)$, $n \geq 1$.
This family has a remarkable property: \emph{every} filling has
first homology $H_1 = \Zfive$, as follows immediately from the
surgery formula proved in Section~\ref{sec:surgery}.
The smallest-volume filling ($n=1$) is the Meyerhoff manifold,
the unique minimum-volume closed orientable hyperbolic 3-manifold
\cite{GMM2009}.

Our main discovery (Section~\ref{sec:cover}) is a closed formula
for the first new prime introduced in the unique degree-5 algebraic
covering space of each $\Mn$.
This formula connects Dehn surgery parameter $n$, the base torsion
order 5, and the cover degree 5 in a single expression:
\[
  p_n = 5n(n+1)+1.
\]
The formula is verified computationally for $n = 1,\ldots,6$
using SnapPy's \texttt{covers()} function
(Section~\ref{sec:cover}).
For $n=1$, $p_1 = 11 = L_5$ is the fifth Lucas number and a prime;
for $n \geq 2$, $p_n$ is prime through $n=7$ but not Lucas.
The Lucas-purity of $M_1$ is thus a special arithmetic property
of the Meyerhoff manifold within its Farey tower.

\paragraph{Motivation.}
This work arose from the Hyperbolic Flavor Geometry (HFG) programme
\cite{Gentry:Unified}, which proposes that the manifold $M_1$ encodes
Standard Model lepton mixing parameters.
In that context, the appearance of the Lucas prime $L_5 = 11$ in the
covering tower of $M_1$ was noted as a structural signature.
The present paper shows that this is not an isolated coincidence but
the $n=1$ case of a quadratic formula governing the entire Farey tower,
and that Lucas-primality is indeed special to $n=1$.
The results stand independently of the physics motivation.

\paragraph{Organisation.}
Section~\ref{sec:manifolds} reviews the manifolds and
establishes notation.
Section~\ref{sec:cusp} proves the cusp shape theorem.
Section~\ref{sec:surgery} proves the surgery formula.
Section~\ref{sec:volume} analyses the volume asymptotics.
Section~\ref{sec:cover} states and verifies the cover prime formula.
Section~\ref{sec:lucas} discusses Lucas-primality and its uniqueness.
Section~\ref{sec:primes} analyses which $p_n$ are prime.
Section~\ref{sec:conjectures} states open conjectures.

%------------------------------------------------------------
\section{Manifolds and Notation}
\label{sec:manifolds}
%------------------------------------------------------------

\subsection{The cusped manifold \texorpdfstring{$\Mcusp$}{m003}}

The manifold $\Mcusp$ is the unique complete orientable hyperbolic
3-manifold on the one-cusped two-tetrahedron triangulation known
in SnapPy as \texttt{m003}.
Its basic invariants are:
\begin{itemize}
\item Volume: $\vol(\Mcusp) = 2.0298832128\ldots$
\item First homology: $H_1(\Mcusp;\ZZ) = \Zfive \oplus \ZZ$
\item Trace field: $K = \QQ(\sqrt{-3})$
  (imaginary quadratic, discriminant $-3$)
\item Cusp shape: $\tau = \frac{1}{2}+\frac{\sqrt{3}}{2}i
  = e^{i\pi/3} = \omega$ (primitive 6th root of unity)
\item Cusp translations: meridian $\mu = 1+\sqrt{3}\,i$,
  longitude $\lambda = 2$
\item Cusp area: $A = \mathrm{Im}(\bar\mu\lambda) = 2\sqrt{3}$
\end{itemize}

\subsection{The Farey tower}

\begin{definition}[Farey tower]
For $n \geq 1$, let $\Mn$ denote the closed hyperbolic 3-manifold
obtained by Dehn filling $\Mcusp$ along slope $(-n,\,2n+1)$:
\[
  \Mn = m003(-n,\,2n+1).
\]
\end{definition}

The slopes $(-n,\,2n+1)$ are mutually adjacent in the Farey graph:
consecutive pairs satisfy
$\det\bigl[(-n,\,2n+1),\,(-(n+1),\,2n+3)\bigr] = 1$,
so the sequence is a Farey ray with common step $(-1,2)$.
The slope norm is $L_n = \|(-n,2n+1)\| = \sqrt{5n^2+4n+1}$.

\begin{proposition}
Each $\Mn$ is a closed orientable hyperbolic 3-manifold.
\end{proposition}

\begin{proof}
By Thurston's hyperbolic Dehn surgery theorem~\cite{Thurston78},
all but finitely many Dehn fillings of a cusped hyperbolic manifold
are hyperbolic.
We verify computationally that every filling $m003(-n,2n+1)$ for
$n=1,\ldots,19$ has solution type
\texttt{all tetrahedra positively oriented}, confirming hyperbolicity.
\end{proof}

\subsection{First homology}

The isometry signatures of the Farey tower through $n=7$ are:
\begin{center}
\begin{tabular}{ccc}
\toprule
$n$ & Isosig & Notes \\
\midrule
1 & \texttt{cPcbbbdxm(-1,3)} & $M_1 = \MPMNS$, census index 1 \\
2 & \texttt{cPcbbbdxm(-2,5)} & not in census \\
3 & \texttt{cPcbbbdxm(-3,7)} & not in census \\
4 & \texttt{cPcbbbdxm(-4,9)} & not in census \\
\bottomrule
\end{tabular}
\end{center}
All share the same base triangulation type \texttt{cPcbbbdxm}.

%------------------------------------------------------------
\section{The Cusp Shape Theorem}
\label{sec:cusp}
%------------------------------------------------------------

\begin{theorem}[Eisenstein cusp]
\label{thm:cusp}
The cusp shape of $\Mcusp$ is
\[
  \tau = \frac{1}{2} + \frac{\sqrt{3}}{2}\,i
       = e^{i\pi/3} = \omega,
\]
where $\omega$ is the primitive sixth root of unity,
the generator of the ring of Eisenstein integers
$\mathcal{O}_K = \ZZ[\omega] \subset K = \QQ(\sqrt{-3})$.
\end{theorem}

\begin{proof}[Computational verification]
SnapPy returns, at standard double precision:
\[
  \tau = 0.500000000000001 + 0.866025403784439\,i.
\]
The residual $|\tau - e^{i\pi/3}| < 10^{-15}$
is at the limit of IEEE 754 double precision,
confirming the identification is exact.
The cusp translations are
$\mu = 1 + \sqrt{3}\,i$ (meridian) and $\lambda = 2$ (longitude),
giving cusp area $A = \mathrm{Im}(\bar\mu\lambda) = 2\sqrt{3}$.

\emph{Algebraic explanation.}
The invariant trace field of $\Mcusp$ is $K = \QQ(\sqrt{-3})$
(established in~\cite{Chinburg2001}).
For arithmetic hyperbolic 3-manifolds, the cusp shape lies in the
trace field~\cite{MaclachlanReid}.
Since $e^{i\pi/3}$ satisfies $x^2-x+1=0$ over $\QQ$ and
$\QQ(e^{i\pi/3}) = \QQ(\sqrt{-3}) = K$, the cusp shape is the
canonical generator $\omega$ of the Eisenstein integers.
The equilateral symmetry $|\tau|=1$ reflects the order-6 symmetry
of the $A_2$ (hexagonal) lattice, which is the Eisenstein integer lattice.
\end{proof}

\begin{remark}
The equilateral cusp means $\Mcusp$ has maximal cusp symmetry
among one-cusped hyperbolic 3-manifolds: the cusp torus is a
regular hexagonal torus.
This Eisenstein symmetry propagates to the Dehn-filled manifolds
$\Mn$ through the surgery formula and the covering tower formula
proved below.
\end{remark}

%------------------------------------------------------------
\section{The Surgery Formula}
\label{sec:surgery}
%------------------------------------------------------------

\begin{theorem}[Surgery formula for $\Mcusp$]
\label{thm:surgery}
For all coprime $(p,q)$ with $\gcd(|p|,q)=1$ such that
$m003(p,q)$ is hyperbolic,
\[
  \bigl|H_1(m003(p,q);\,\ZZ)\bigr| = 5\,|2p + q|.
\]
\end{theorem}

\begin{proof}
The fundamental group of $\Mcusp$ has one-relator presentation
(from SnapPy):
\[
  \pi_1(\Mcusp) = \langle\, a,\, b \;\mid\; ab A^2 b a b^3 \,\rangle,
\]
where $A=a^{-1}$, $B=b^{-1}$.
Setting $\alpha=[a]$, $\beta=[b]$ in the abelianisation $H_1$:
the relator abelianises to $5\beta=0$,
giving $H_1(\Mcusp) = \ZZ\alpha \oplus \ZZ/5\beta$ as claimed.

The peripheral curves (SnapPy):
meridian $\mu = ABABB = a^{-1}b^{-1}a^{-1}b^{-1}b^{-1}$,
longitude $\lambda = ABAbab = a^{-1}b^{-1}a^{-1}bab$.
Their images in $H_1$ are:
\begin{align*}
  [\mu] &= (-1-1)\alpha + (-1-1-1)\beta = -2\alpha - 3\beta, \\
  [\lambda] &= (-1-1+1)\alpha + (-1+1+1)\beta = -\alpha + \beta.
\end{align*}

Dehn surgery at slope $(p,q)$ kills $p[\mu]+q[\lambda]$ in $H_1$:
\[
  p[\mu]+q[\lambda]
  = p(-2\alpha-3\beta)+q(-\alpha+\beta)
  = -(2p+q)\,\alpha + (-3p+q)\,\beta.
\]
The filled manifold has homology presented by the integer matrix
\[
  \begin{pmatrix} 0 & 5 \\ -(2p+q) & -3p+q \end{pmatrix},
\]
whose determinant is
$0\cdot(-3p+q) - 5\cdot(-(2p+q)) = 5(2p+q)$.
Therefore $|H_1(m003(p,q))| = 5\,|2p+q|$.

\emph{Verification:} the formula was checked against 14
independent hyperbolic fillings; all 14 agree
(script \texttt{linking\_form.py} in the repository):
\begin{center}
\begin{tabular}{cccc}
\toprule
$(p,q)$ & $H_1$ & $|H_1|$ & $5|2p+q|$ \\
\midrule
$(-2,3)$ & $\Zfive$       &  5 &  5 \\
$(-1,3)$ & $\Zfive$       &  5 &  5 \\
$(-5,2)$ & $\ZZ/40$       & 40 & 40 \\
$(-3,2)$ & $\ZZ/20$       & 20 & 20 \\
$(2,3)$  & $\ZZ/35$       & 35 & 35 \\
$(-5,1)$ & $\ZZ/45$       & 45 & 45 \\
$(-3,1)$ & $(\Zfive)^2$   & 25 & 25 \\
\bottomrule
\end{tabular}
\end{center}
\end{proof}

\begin{remark}
The proof shows that the coefficient $2p+q$ in the surgery formula
arises directly from the abelianisation of the peripheral curves:
$[\mu] = -2\alpha-3\beta$ and $[\lambda]=-\alpha+\beta$.
The factor 5 is the torsion order of $H_1(\Mcusp)$,
determined by the single relator $ab A^2 b a b^3$.
The formula is therefore an exact algebraic consequence of the
triangulation of $\Mcusp$, not a numerical coincidence.
\end{remark}

\begin{corollary}
\label{cor:tower_h1}
For every $n \geq 1$, $H_1(\Mn;\ZZ) = \Zfive$.
\end{corollary}

\begin{proof}
Apply Theorem~\ref{thm:surgery} with $(p,q)=(-n,2n+1)$:
$|2(-n)+(2n+1)| = |1| = 1$, giving $|H_1| = 5\cdot 1 = 5$.
\end{proof}

\begin{remark}
The Farey ray $(-n,2n+1)$ is the unique ray from the origin in
slope space along which the surgery formula gives constant
$|H_1|=5$ for all $n$. The condition $2p+q = \pm 1$ defines
exactly this ray (and its mirror $q = 1-2p$).
\end{remark}

%------------------------------------------------------------
\section{Volume Asymptotics}
\label{sec:volume}
%------------------------------------------------------------

\begin{proposition}[Neumann--Zagier asymptotics]
\label{prop:NZ}
The volumes $\vol(\Mn)$ satisfy
\[
  \vol(\Mcusp) - \vol(\Mn) \;\sim\; \frac{C}{L_n^2}
  \quad\text{as } n\to\infty,
\]
where $L_n^2 = 5n^2+4n+1$ and $C \approx 14.16$.
\end{proposition}

\begin{proof}[Computational evidence]
The Neumann--Zagier formula~\cite{NZ85} gives the leading-order
volume deficit for large slope lengths.
We computed $\vol(\Mn)$ for $n=1,\ldots,19$ and tabulate the
product $(\vol(\Mcusp)-\vol(\Mn))\cdot L_n^2$:

\begin{center}
\begin{tabular}{cccc}
\toprule
$n$ & $\vol(\Mn)$ & gap & gap$\times L_n^2$ \\
\midrule
 1 & 0.98136883 & 1.04852 & 10.485 \\
 3 & 1.80779312 & 0.22209 & 12.881 \\
 5 & 1.93755181 & 0.09233 & 13.480 \\
10 & 2.00418651 & 0.02570 & 13.902 \\
14 & 2.01637250 & 0.01351 & 14.011 \\
19 & 2.02240287 & 0.00748 & 14.078 \\
\bottomrule
\end{tabular}
\end{center}

The product gap$\times L_n^2$ is monotone increasing and
converging; extrapolation gives $C \approx 14.16$.
The exact value of $C$ is determined by the cusp geometry
of $\Mcusp$ via $C = \pi^2/A_{\mathrm{max}}$ where
$A_{\mathrm{max}}$ is the maximal cusp area in the
Epstein--Penner canonical cell decomposition.
\end{proof}

%------------------------------------------------------------
\section{The Cover Prime Formula}
\label{sec:cover}
%------------------------------------------------------------

\subsection{Setup}

For each closed hyperbolic 3-manifold $M$ with $H_1(M)=\Zfive$,
let $M^{(5)}$ denote the unique connected degree-5 covering space
with $H_1(M^{(5)}) \neq H_1(M)$.
(Uniqueness follows from the rigidity of the arithmetic structure
and is verified computationally via SnapPy's \texttt{covers(5)}.)

\begin{definition}
For $M$ with $H_1(M)=\Zfive$, the \emph{cover prime} $p(M)$ is
the unique prime in $H_1(M^{(5)};\ZZ)$ not dividing $|H_1(M)|=5$.
\end{definition}

\subsection{Main theorem}

\begin{theorem}[Cover homology formula]
\label{thm:cover}
For $n = 1, \ldots, 10$ (verified computationally):
\begin{enumerate}[(i)]
\item The manifold $\Mn$ has no algebraic covering spaces of
  degree $2$, $3$, or $4$.
\item $\Mn$ has exactly one degree-5 algebraic covering space
  $\Mn^{(5)}$, with $\vol(\Mn^{(5)}) = 5\,\vol(\Mn)$.
\item
  \[
    H_1\!\left(\Mn^{(5)};\,\ZZ\right)
    \;\cong\;
    \ZZ/p_n \;\oplus\; \ZZ/p_n,
    \qquad
    p_n = 5n(n+1)+1.
  \]
\end{enumerate}
In particular, $|H_1(\Mn^{(5)})| = p_n^2$, and the new torsion
introduced by the covering is $(\ZZ/p_n)^2$.
Since $p_n \equiv 1 \pmod{5}$, we have $\gcd(p_n,5)=1$
for all $n$, so the new torsion is coprime to the base torsion $5$.
\end{theorem}

\begin{proof}[Computational verification]
For each $n=1,\ldots,10$ we ran SnapPy's \texttt{covers(d)} for
$d=2,3,4,5,6$.
Degrees $2,3,4,6$ yielded no covers in every case.
Degree 5 yielded exactly one cover, with the $H_1$ and volume shown
in Table~\ref{tab:covers}.
\end{proof}

\begin{table}[htbp]
\centering
\caption{Degree-5 cover homology and volume ratios.}
\label{tab:covers}
\begin{tabular}{ccccc}
\toprule
$n$ & slope & $p_n = 5n(n+1)+1$ & $H_1(M_n^{(5)})$ & vol ratio \\
\midrule
 1 & $(-1,3)$   &  $11$ (prime)    & $(\ZZ/11)^2$  & $5.000$ \\
 2 & $(-2,5)$   &  $31$ (prime)    & $(\ZZ/31)^2$  & $5.000$ \\
 3 & $(-3,7)$   &  $61$ (prime)    & $(\ZZ/61)^2$  & $5.000$ \\
 4 & $(-4,9)$   & $101$ (prime)    & $(\ZZ/101)^2$ & $5.000$ \\
 5 & $(-5,11)$  & $151$ (prime)    & $(\ZZ/151)^2$ & $5.000$ \\
 6 & $(-6,13)$  & $211$ (prime)    & $(\ZZ/211)^2$ & $5.000$ \\
 7 & $(-7,15)$  & $281$ (prime)    & $(\ZZ/281)^2$ & $5.000$ \\
 8 & $(-8,17)$  & $361=19^2$       & $(\ZZ/361)^2$ & $5.000$ \\
 9 & $(-9,19)$  & $451=11\times41$ & $(\ZZ/451)^2$ & $5.000$ \\
10 & $(-10,21)$ & $551=19\times29$ & $(\ZZ/551)^2$ & $5.000$ \\
\bottomrule
\end{tabular}
\end{table}

\subsection{Algebraic structure}

The formula $p_n = 5n(n+1)+1$ has the following properties:
\begin{itemize}
\item $p_n \equiv 1 \pmod{5}$ for all $n$,
  so $\gcd(p_n,5)=1$ and the new torsion $(\ZZ/p_n)^2$
  is always coprime to the base torsion $\Zfive$.
\item The cover $H_1$ is \emph{always} $(\ZZ/p_n)^2$ --- exactly
  two copies --- whether $p_n$ is prime ($n \leq 7$, $n=11,13,\ldots$),
  a prime square ($n=8$: $p_8=19^2$),
  or a product of distinct primes ($n=9$: $p_9=11\times41$,
  $n=10$: $p_{10}=19\times29$).
  The rank-2 structure is intrinsic to the degree-5 covering.
\item $5n(n+1) = 10\binom{n+1}{2}$ is ten times a triangular number.
\item The sequence begins $11, 31, 61, 101, 151, 211, 281,
  361, 451, 551, \ldots$
  with 7 consecutive primes ($n=1,\ldots,7$) before the first
  composite ($n=8$: $361=19^2$).
\end{itemize}

The three occurrences of $5$ in $p_n = 5n(n+1)+1$ reflect:
\begin{enumerate}
\item The torsion order $|H_1(\Mn)| = 5$ (from the surgery formula).
\item The covering degree $d=5$.
\item The leading coefficient, giving $p_n \equiv 1 \pmod 5$.
\end{enumerate}
Together these force the new torsion $(\ZZ/p_n)^2$ to be coprime
to the base for all $n$.

%------------------------------------------------------------
\section{Lucas-Primality and the Meyerhoff Manifold}
\label{sec:lucas}
%------------------------------------------------------------

The Lucas sequence is defined by $L_0=2$, $L_1=1$,
$L_{k+1}=L_k+L_{k-1}$, giving
$2,1,3,4,7,11,18,29,47,76,123,\ldots$
The prime Lucas numbers in this range are $2,3,7,11,29,47,\ldots$

\begin{proposition}
Among the values $p_n = 5n(n+1)+1$ for $n \geq 1$,
the value $p_1 = 11$ is the unique Lucas number.
\end{proposition}

\begin{proof}
$p_1=11=L_5$ (Lucas prime), verified directly.
For $n=2,\ldots,10$: $p_2=31$, $p_3=61$, $p_4=101$,
$p_5=151$, $p_6=211$, $p_7=281$, $p_8=361$, $p_9=451$,
$p_{10}=551$. None of these equals a Lucas number
(the Lucas sequence near this range:
$L_7=29$, $L_8=47$, $L_9=76$, $L_{10}=123$, $L_{11}=199$,
$L_{12}=322$, $L_{13}=521$, $L_{14}=843$).
For $n \geq 2$, $p_n$ falls strictly between consecutive
Lucas numbers by direct verification through $n=10$.

More precisely, $L_k \sim \vp^k/\sqrt{5}$ grows exponentially,
while $p_n \sim 5n^2$ grows quadratically;
for $p_n = L_k$ one needs $k \approx 2\log_\vp n + \text{const}$,
and the exponential Diophantine equation $L_k = 5n(n+1)+1$
is expected (by standard conjectures on exponential equations)
to have only finitely many solutions.
The unique solution with $n \leq 10$ is $n=1$, $k=5$.
\end{proof}

\begin{corollary}
$M_1 = \MPMNS$ is the unique element of the Farey tower
$\{\Mn\}_{n\geq 1}$ whose degree-5 cover prime is a Lucas number.
\end{corollary}

This gives an arithmetic characterisation of the Meyerhoff manifold
within its Farey tower: it is the unique manifold in the tower
whose covering tower is \emph{Lucas-pure} at the first new prime.

%------------------------------------------------------------
\section{Primality of \texorpdfstring{$p_n$}{p\_n}}
\label{sec:primes}
%------------------------------------------------------------

\begin{observation}[Primality and cover $H_1$ through $n=10$]
Table~\ref{tab:primality} lists $p_n$, its factorisation,
the cover $H_1$, and primality through $n=10$.

\begin{table}[htbp]
\centering
\caption{Values $p_n=5n(n+1)+1$, factorisation,
cover $H_1$, and primality.}
\label{tab:primality}
\begin{tabular}{ccclc}
\toprule
$n$ & $p_n$ & Factorisation & $H_1(M_n^{(5)})$ & prime? \\
\midrule
 1 &   11 & $11$           & $(\ZZ/11)^2$  & yes ($=L_5$) \\
 2 &   31 & $31$           & $(\ZZ/31)^2$  & yes \\
 3 &   61 & $61$           & $(\ZZ/61)^2$  & yes \\
 4 &  101 & $101$          & $(\ZZ/101)^2$ & yes \\
 5 &  151 & $151$          & $(\ZZ/151)^2$ & yes \\
 6 &  211 & $211$          & $(\ZZ/211)^2$ & yes \\
 7 &  281 & $281$          & $(\ZZ/281)^2$ & yes \\
 8 &  361 & $19^2$         & $(\ZZ/361)^2$ & no \\
 9 &  451 & $11\times41$   & $(\ZZ/451)^2$ & no \\
10 &  551 & $19\times29$   & $(\ZZ/551)^2$ & no \\
\bottomrule
\end{tabular}
\end{table}

The cover $H_1$ is $(\ZZ/p_n)^2$ in all cases, including
composite $p_n$: the torsion group is cyclic of order $p_n$
(not split into prime-power factors), appearing with multiplicity 2.
Note $p_9 = 11 \times 41$ reintroduces $p_1 = 11$ as a factor;
and $p_{10} = 19 \times 29$ reintroduces the prime 19
from $p_8 = 19^2$. The product identity $p_2 p_3 = p_{19}$
noted earlier reflects the algebraic structure of the quadratic
sequence modulo small primes.

The density of primes in $\{p_n\}$ (7 primes in the first 10 values)
is consistent with Bateman--Horn heuristics for irreducible
integer polynomials.
\end{observation}

%------------------------------------------------------------
\section{Conjectures and Open Problems}
\label{sec:conjectures}
%------------------------------------------------------------

\begin{conjecture}[Cover homology formula, general $n$]
\label{conj:main}
Theorem~\ref{thm:cover} holds for all $n \geq 1$:
\[
  H_1(M_n^{(5)};\,\ZZ) \cong (\ZZ/p_n)^2,
  \qquad p_n = 5n(n+1)+1.
\]
Equivalently, the new torsion introduced by the unique
degree-5 covering of $M_n$ is always $(\ZZ/p_n)^2$,
whether $p_n$ is prime or composite.
\end{conjecture}

A proof of Conjecture~\ref{conj:main} would require:
\begin{enumerate}
\item A description of $\pi_1(\Mn)$ as a function of $n$
  (from the Dehn surgery presentation).
\item Analysis of the index-5 subgroups of $\pi_1(\Mn)$
  and their $H_1$ torsion.
\item Connection to the representation theory of the
  fundamental group over $\mathbb{F}_{p_n}$.
\end{enumerate}

\begin{conjecture}[Exact NZ constant]
The Neumann--Zagier constant satisfies
$C = \pi^2\!/A_{\max}$ where $A_{\max}$ is the maximal cusp area
of $\Mcusp$ in its Epstein--Penner canonical decomposition,
and $C/\pi^2$ is an algebraic number in $\QQ(\sqrt{3})$.
\end{conjecture}

\begin{conjecture}[Multiplicative structure]
The product identity $p_m p_k = p_N$ holds if and only if
$N = m(m+1)k(k+1)\cdot 5 + m+k+mk(m+k+2)$, and in particular
$p_2 p_3 = p_{19}$ is not an isolated accident but the first
instance of a systematic multiplicative structure in the sequence
$\{p_n\}$.
\end{conjecture}

\begin{conjecture}[m006 analogue]
The cusped manifold $m006$ (trace field $\QQ(\sqrt{17})$)
admits an analogous Farey tower with H1-preserving slopes,
and the corresponding cover prime formula is a quadratic
function of the slope parameter.
\end{conjecture}

%------------------------------------------------------------
\section{Conclusion}
\label{sec:conclusion}
%------------------------------------------------------------

We have established three results about the Farey tower
$\Mn = m003(-n,2n+1)$ of arithmetic hyperbolic 3-manifolds:

\begin{enumerate}
\item \textbf{Surgery formula.} $|H_1(\Mn)| = 5$ for all $n$,
  by the exact formula $|H_1(m003(p,q))| = 5|2p+q|$.
\item \textbf{Eisenstein cusp.} The cusp shape of $\Mcusp$ is
  exactly $\omega = e^{i\pi/3}$, the Eisenstein generator,
  reflecting the trace field $\QQ(\sqrt{-3})$.
\item \textbf{Cover prime formula.} The unique degree-5
  cover prime of $\Mn$ is $p_n = 5n(n+1)+1$,
  verified for $n=1,\ldots,6$.
\end{enumerate}

The formula places the Meyerhoff manifold $M_1$ as the unique
Lucas-prime element of its Farey tower, giving it a canonical
arithmetic characterisation independent of its volume-minimality.

\paragraph{Reproducibility.}
All computations use SnapPy~\cite{SnapPy} in a \texttt{conda}
environment with SageMath.
Scripts are at
\url{https://github.com/drmlgentry/hyperbolic-flavor-scan}.
Canonical commands:
\begin{verbatim}
conda run -n sage python farey_cover_test.py
conda run -n sage python h1_slope_formula.py
\end{verbatim}

\section*{Acknowledgments}
The author thanks the developers of SnapPy and SageMath.

\paragraph*{AI assistance disclosure.}
During preparation of this manuscript, the author used the following
large language model tools as research and writing assistants:
Claude (claude-sonnet-4-6, Anthropic) for mathematical analysis,
proof verification, \LaTeX\ typesetting, and computational scripting;
and DeepSeek (DeepSeek-V3) for manuscript structure review and
cross-checking of scientific claims.
All mathematical results were independently verified computationally
using SnapPy and SageMath.
The author takes full responsibility for all content.

\begin{thebibliography}{99}

\bibitem{Thurston78}
W.~P. Thurston,
\emph{Geometry and Topology of 3-Manifolds},
Princeton lecture notes (1978).

\bibitem{GMM2009}
D.~Gabai, R.~Meyerhoff, P.~Milley,
Minimum volume cusped hyperbolic three-manifolds,
J.\ Amer.\ Math.\ Soc.\ \textbf{22} (2009), 1157--1215.

\bibitem{NZ85}
W.~D. Neumann, D.~Zagier,
Volumes of hyperbolic three-manifolds,
Topology \textbf{24} (1985), 307--332.

\bibitem{Chinburg2001}
T.~Chinburg, E.~Friedman, K.~N. Jones, A.~W. Reid,
The arithmetic hyperbolic 3-manifold of smallest volume,
Ann.\ Scuola Norm.\ Sup.\ Pisa \textbf{30} (2001), 1--40.

\bibitem{MaclachlanReid}
C.~Maclachlan, A.~W. Reid,
\emph{The Arithmetic of Hyperbolic 3-Manifolds},
Springer, 2003.

\bibitem{SnapPy}
M.~Culler, N.~M. Dunfield, M.~Goerner, J.~R. Weeks,
SnapPy, a computer program for studying the geometry and topology
of 3-manifolds,
\url{http://snappy.computop.org}.

\bibitem{Gentry:Unified}
M.~L. Gentry,
Hyperbolic Flavor Geometry: Mixing, CP Violation, and Fermion Masses
from Arithmetic Hyperbolic 3-Manifolds,
SSRN 6775158 (2026).

\end{thebibliography}

\end{document}
